\section{Reading only what the algorithm requires}
\label{sec:lazy-sequence-facade}
\label{sec:sequence-materialization}

The storage can be provided in different forms as explained previously in the first section~\ref{sec:sequence-abstractions}
In BioseqDB main implementation of Postgres-based storage is done via \path{lazy_sequence_view<A>} and \path{retained_lazy_sequence_view<A>} mentioned in table~\ref{tab:sequence-representation-specializations}.
These types holds pointers to functions that implement the storage and represent the following operations:
\begin{itemize}
  \item Report sizes of sequences
  \item Copy selected intervals into caller-owned memory, which in this case can be GPU-owned memory in GPU-accelerated pipeline stages
  \item Reading successive chunks
  \item Checking if the compatible interval was already materialized and can be borrowed (caching)
\end{itemize}
Internally I distinguish several access policies which are specified by the pipeline stage, they are the following choices:
\begin{enumerate}
  \item \textbf{Stream.} If the algorithm accepts a stream, read the selected
        interval in bounded chunks and save it to reusable scratch memory (large memory pool)
  \item \textbf{Borrow.} If the algorithm needs packed data, first check if's already in cache and safe to borrow.
  \item \textbf{Pack.} In case if borrowing is not possible, algorithm needs to request new memory permit (this is routed through custom pool allocator). Caller is then responsible for this memory chunk.
\end{enumerate}

\begin{figure}[H]
\centering
\begin{tikzpicture}[x=1cm,y=1cm,font=\footnotesize,
  request/.style={draw=black!60,fill=black!3,minimum width=3.9cm,
    minimum height=.72cm,align=center},
  decision/.style={diamond,aspect=2.7,draw=black!60,
    inner sep=3pt,align=center}]
  \node[request] (q) at (6.5,4.6) {requested interval $[a,b)$};
  \node[decision] (need) at (6.5,3.25) {one packed block needed?};
  \draw[seq nine flow] (q)--(need);
  \coordinate (stream) at (1.6,1.05);
  \draw[seq nine flow,seqNineBlue] (need.west)-|node[pos=.18,above] {no} (stream);
  \node[decision] (safe) at (9.05,1.9) {safe to borrow?};
  \draw[seq nine flow] (need.east)-|node[pos=.3,above] {yes} (safe);
  \draw[seq nine flow,seqNineGreen] (safe.west)-|node[pos=.27,above] {yes} (6.45,.45);
  \draw[seq nine flow,seqNineOrange] (safe.east)-|node[pos=.27,above] {no} (11.5,.45);
  \foreach \r in {0,1,2} {
    \foreach \i in {0,1,2}
      \draw[draw=seqNineBlue,fill=seqNineBlue!12]
        (.5+.57*\i,.3-.5*\r) rectangle ++(.52,.36);
  }
  \draw[seq nine flow,seqNineBlue] (2.35,.46)--(2.35,-.54);
  \node[align=center] at (1.55,-1.35)
    {stream chunks\\reuse scratch memory};
  \draw[fill=black!4,draw=black!50] (4.4,-.32) rectangle (8.3,.2);
  \draw[draw=seqNineGreen,fill=seqNineGreen!20,thick]
    (5.375,-.32) rectangle (7.325,.2);
  \foreach \x in {5.375,6.35,7.325}
    \draw[black!55] (\x,-.32)--(\x,.2);
  \node[below,text=seqNineGreen] at (6.35,-.4) {borrow this interval};
  \node[align=center] at (6.35,-1.5)
    {existing packed units\\owner stays alive};
  \foreach \i in {0,...,4}
    \draw[seqNineOrange,fill=seqNineOrange!15]
      (10.05+.52*\i,-.32) rectangle ++(.52,.52);
  \node[below,text=seqNineOrange] at (11.35,-.4) {new packed interval};
  \node[align=center] at (11.35,-1.5)
    {allocate with a permit\\retain buffer and permit};
\end{tikzpicture}
\caption{Simplified view of different access policies}
\label{fig:sequence-abstraction-layers}
\end{figure}

\info{Borrowing}{Borrowing can fail if the region for example starts in the middle of storage unit in packed sequence, then we cannot reuse that memory chunk as it's misaligned}

The following pseudocode states the logical read rule. A piece may refer to
memory or to a stored fragment. Locating it and fetching its storage are
handled by the storage reader.
\begin{algorithm}[H]
\caption{Read one interval of a composed sequence}
\label{lst:lazy-window-read}
\small
\begin{algorithmic}[1]
  \Require view $V$, valid bounds $0\leq a\leq b\leq |V|$
  \Require output space for $b-a$ ranks
  \ForAll{pieces intersecting $V[a,b)$, in logical order}
    \State Map the intersection to that piece's source positions
    \State Read those positions, using bounded scratch space if needed
    \State Apply the piece's reverse and complement flags
    \State Append the resulting ranks to the output
  \EndFor
  \State Require exactly $b-a$ output ranks
\end{algorithmic}
\end{algorithm}

An interval crossing two fragments must still produce the same symbols as an
equivalent flat value. A storage reader may fetch complete fragments, but
only the requested ranks enter the logical output. Thus the interface limits
what the algorithm asks for without claiming that every source supports
equally small physical reads.

\subsection{Where a flat value is required}
\label{sec:sequence-stdc}

BioSeqDB's flat PostgreSQL value consists of a fixed header followed by aligned
packed data. The header records the alphabet, unit width, length, and padding.
The adapter validates the header and payload before selecting a typed compact
range. Runtime checks happen at this boundary, so the inner loop can use
compile-time packing rules.

An already resident, uncompressed, aligned value can be borrowed through this
adapter. PostgreSQL may first have to fetch out-of-line data or decompress it
using TOAST.\footcite{postgresql17_toast} That work precedes borrowing and can
include a copy.

Writing a lazy expression into durable storage is another boundary. Its
temporary source references cannot be stored as raw pointers. The expression
must be flattened to a value that can stand on its own. Text output similarly
has to create the requested text. If the result exceeds the flat PostgreSQL
value's size limit, the operation reports an error.

For repeated full reads, one packed copy can be a useful choice. Laziness
preserves that option while also supporting the small-interval case. It is
the consumer's access requirement, not the mere existence of a slice, that
determines when to build the copy.
